% Estilo comum do relatório: código Lastro, código C e saída do compilador.
% Em tons de cinza, como recomenda o modelo SBC.

% Código curto no meio do texto e nas tabelas.
\newcommand{\cod}[1]{\texttt{#1}}

% Sem esticar o espaço entre parágrafos para encher a página: com tabelas e listagens,
% o \flushbottom do modelo abria buracos antes dos títulos.
\raggedbottom

\makeatletter
% O modelo só define \section e \subsection; este nível segue o mesmo padrão, em itálico.
\renewcommand\subsubsection{\@startsection{subsubsection}{3}{\z@}%
                       {-6\p@ \@plus -2\p@ \@minus -2\p@}%
                       {0.1\p@}%
                       {\normalfont\normalsize\bfseries\itshape}}

% O hyperref troca o \@lbibitem do modelo e volta a imprimir o rótulo [Autor Ano] na
% margem. Refaz a versão do modelo (referência sem rótulo), mantendo a âncora do link.
\def\@lbibitem[#1]#2{\item[]%
  \Hy@raisedlink{\hyper@anchorstart{cite.#2\@extra@b@citeb}\hyper@anchorend}%
  \if@filesw{\let\protect\noexpand\immediate
    \write\@auxout{\string\bibcite{#2}{#1}}}\fi\ignorespaces}
\makeatother

\lstdefinelanguage{Lastro}{
  morekeywords={const,proc,func,int,float,bool,if,else,while,for,to,
                break,continue,return,ref,true,false},
  sensitive=true,
  morecomment=[l]{//},
  morestring=[b]",
}

\lstset{
  basicstyle=\ttfamily\small,
  keywordstyle=\bfseries,
  commentstyle=\itshape\color{black!60},
  stringstyle=\color{black!75},
  columns=fullflexible,
  keepspaces=true,
  showstringspaces=false,
  upquote=true,
  frame=lines,
  framerule=0.4pt,
  rulecolor=\color{black!40},
  framesep=4pt,
  xleftmargin=0.8cm,
  xrightmargin=0.8cm,
  aboveskip=6pt,
  belowskip=6pt,
  % acentos em comentários e textos (listings lê byte a byte)
  literate={á}{{\'a}}1 {à}{{\`a}}1 {â}{{\^a}}1 {ã}{{\~a}}1
           {é}{{\'e}}1 {ê}{{\^e}}1 {í}{{\'i}}1 {ó}{{\'o}}1
           {ô}{{\^o}}1 {õ}{{\~o}}1 {ú}{{\'u}}1 {ç}{{\c{c}}}1
           {Á}{{\'A}}1 {É}{{\'E}}1 {Ó}{{\'O}}1 {Ç}{{\c{C}}}1,
}

\lstdefinestyle{lastro}{language=Lastro}
\lstdefinestyle{c}{language=C}
% saída do terminal: sem destaque de palavras
\lstdefinestyle{saida}{language={}, basicstyle=\ttfamily\small\color{black!85}}
